\section{Conclusion}
\label{sec:conclusion}

Given an explicit channel for doubts, tool-using scientific agents from five vendors use it---flagging three-residue proteins, bacterial proteins with synaptic localisations, molecules of two daltons and negative donor counts---and then answer with the value anyway unless the record gives them an exact route to a replacement, or an instruction tells them to check. Plausible corruptions pass through almost regardless; in chemistry they are often detected and pass through anyway. In our tasks the detection that does occur behaves like a consistency check and a memory lookup rather than a plausibility judgement, and the correction that does occur is a computation the record made possible. The practical lesson is that records returned to agents should be designed to be checkable, and harnesses should be designed to act on what the agent says it noticed.

\section*{Impact Statement}
This work studies how language-model agents handle corrupted database records in scientific question answering. Its direct effect is on the design of agent harnesses and tool interfaces for the life sciences and chemistry, where silently propagated errors could reach experimental decisions; we release the benchmark so that such failures can be measured before deployment. The corruptions are synthetic and the records public. We see no societal risk specific to this work beyond those general to evaluating and deploying language-model agents.
